\subsection{Arbitrary weighted sums}
\label{sec:big-unrestricted-addition}

With existence for general measures available, weighted addition extends to arbitrary $u,v\in\mathcal E_\varphi$. Truncating the exponential weights gives finite data, which \cref{thm:big-model-exponential} solves. A first-moment estimate prevents collapse of these solutions, and a rooftop involving their limit localizes the equation before the limit can be compared with the inputs.

\Needspace{11\baselineskip}
\begin{theorem}\label{thm:big-unrestricted-addition}\index{Monge--Ampere measure@Monge--Amp\`ere measure!addition in relative classes}
For every $u,v\in\mathcal E_\varphi$ and $\lambda>0$, there is a unique $a\in\mathcal E_\varphi$ satisfying
\[
 \mu_a=\mathrm{e}^{\lambda(a-u)}\mu_u+\mathrm{e}^{\lambda(a-v)}\mu_v.
\]
Moreover $a\leq\min(u,v)$ and $\mu_a\leq\mu_u+\mu_v$.
\end{theorem}
\begin{proof}
By the scaling and the common translation used in \cref{thm:big-weighted-addition}, it suffices to take $\lambda=1$ and $u,v\leq0$. For positive integers $j$, let
\[
 \tau_j=\mathrm{e}^{-\max(u,-j)}\mu_u+\mathrm{e}^{-\max(v,-j)}\mu_v.
\]
These are nonzero finite non-pluripolar measures, and they increase with $j$. Only the weights are truncated; the functions $\max(u,-j)$ and $\max(v,-j)$ need not be $\theta$-plurisubharmonic. \Cref{thm:big-model-exponential} gives $a_j\in\mathcal E_\varphi$ with $\mu_{a_j}=\mathrm{e}^{a_j}\tau_j$. Since $\tau_j$ increases, exponential order implies that $a_j$ decreases.

Set $r=P_\theta(u,v)\in\mathcal E_\varphi$. On $\{a_j<r\}$, one has $r\leq u,v$ and
\[
 (r-a_j)\mathrm{e}^{a_j-\max(u,-j)}\leq(r-a_j)\mathrm{e}^{a_j-r}\leq \mathrm{e}^{-1},
\]
with the same inequality for $v$. Consequently
\[
 \int_X(r-a_j)_+\,\mu_{a_j}
 \leq \mathrm{e}^{-1}\bigl(\mu_u(X)+\mu_v(X)\bigr).
\]
Also $a_j\leq a_1\leq\varphi+\sup_Xa_1$. The one-sided first-moment criterion of \cref{prop:big-relative-first-moment}, with fixed reference $r$, bounds $\sup_Xa_j$ from below. Thus the decreasing sequence has an $L^1$ limit $a\not\equiv-\infty$, and the same criterion puts $a$ in $\mathcal E_\varphi$.

We identify the equation before comparing $a$ with the inputs. The rooftop
\[
 q=P_\theta(a,u,v)\in\mathcal E_\varphi
\]
lies below $a_j$, $a$, $u$ and $v$.

\Needspace{10\baselineskip}
For $T>0$, let $\chi_T=g_T(q)$, with $g_T$ as in the proof of \cref{lem:big-truncated-addition-limit}: a continuous nondecreasing cutoff that vanishes where $q\leq-2T$, equals one where $q\geq-T$, and increases to one off the pole set of $q$ as $T\to\infty$. On $\{\chi_T>0\}$ both inputs exceed $-2T$. Hence, for $j>2T$,
\begin{equation}\label{eq:unrestricted-addition-cutoff}
 \chi_T\mu_{a_j}
 =\chi_T \mathrm{e}^{a_j}(\mathrm{e}^{-u}\mu_u+\mathrm{e}^{-v}\mu_v).
\end{equation}
The factors $\chi_T\mathrm{e}^{-u}$ and $\chi_T\mathrm{e}^{-v}$ are bounded by $\mathrm{e}^{2T}$. Since $a_j\downarrow a$ and $a_j\leq\sup_Xa_1$, dominated convergence against the fixed measure $\mu_u+\mu_v$ gives convergence in total variation of the right-hand sides to the corresponding expression with $a$.

For the left-hand side, $a_j\to a$ in capacity by \cref{lem:monotone-psh-capacity}. Choose a Hermitian form $\Omega\geq\theta$, and replace $a_j$ by $b_j=\max(a_j,-2T)$ before forming bounded products. The $b_j$ are uniformly bounded $\Omega$-plurisubharmonic functions converging in capacity. The expansion
\[
 (\theta+\ddc b_j)^n
 =\sum_{k=0}^n\binom nk(\Omega+\ddc b_j)^k
                         \wedge(\theta-\Omega)^{n-k}
\]
and bounded capacity continuity apply against the fixed quasi-continuous weight $\chi_T$. Where $\chi_T\ne0$ no potential is truncated, so locality gives convergence to $\chi_T\mu_a$. In this way the possibly signed bounded product is handled before the cutoff is inserted.

Thus \eqref{eq:unrestricted-addition-cutoff} holds in the limit. Let $T\to\infty$. Both sides ignore the poles of $q$, so monotone convergence gives
\[
 \mu_a=\mathrm{e}^{a-u}\mu_u+\mathrm{e}^{a-v}\mu_v.
\]
Since $a\in\mathcal E_\varphi$, this also shows that the weighted right-hand side is finite.

For $\varepsilon>0$, the equation implies
\[
 \mathbf1_{\{u<a-\varepsilon\}}\mu_u
 \leq \mathrm{e}^{-\varepsilon}\mathbf1_{\{u<a-\varepsilon\}}\mu_a.
\]
Strict domination, \cref{thm:big-relative-domination}, gives $u\geq a-\varepsilon$, and letting $\varepsilon\downarrow0$ gives $u\geq a$. Similarly $v\geq a$, so $\mu_a\leq\mu_u+\mu_v$. Two solutions have the same exponentially weighted measure, so \cref{cor:big-relative-orders}, applied in both directions, gives uniqueness.
\end{proof}

\begin{corollary}\label{cor:weighted-addition-algebra}
For fixed $\lambda>0$, write $u\star_\lambda v$ for the solution $a$ in the theorem. This operation on $\mathcal E_\varphi$ is commutative and associative, and, for $c\in\mathbb R$,
\[
 (u+c)\star_\lambda(v+c)=(u\star_\lambda v)+c,
 \qquad
 u\star_\lambda u=u-\lambda^{-1}\log2.
\]
\end{corollary}
\begin{proof}
Put $F_\lambda(u)=\mathrm{e}^{-\lambda u}\mu_u$, with zero mass on the poles. The theorem says that
$F_\lambda(u\star_\lambda v)=F_\lambda(u)+F_\lambda(v)$.
By \cref{cor:big-relative-orders}, the map $F_\lambda$ is injective on $\mathcal E_\varphi$. Commutativity and associativity therefore follow from those of addition of positive measures, including measures of infinite total mass. Finally,
$F_\lambda(u+c)=\mathrm{e}^{-\lambda c}F_\lambda(u)$, which together with injectivity gives both identities.
\end{proof}
